\documentclass[10pt,a4paper]{amsart}
\usepackage[margin=19mm]{geometry}
\usepackage{amsmath,amssymb,mathtools}
\usepackage[scr=rsfs,cal=euler]{mathalfa}
\usepackage{xfrac}
\usepackage[unicode,hidelinks,bookmarks=false]{hyperref}
\newcommand{\ie}{i.e.\ }
\newcommand{\cf}{cf.\ }
\newcommand{\resp}{respectively\ }
\newcommand{\Subsection}{\subsection}
\providecommand{\nobreakdash}{\nobreak\hbox{-}}
\setlength{\emergencystretch}{2em}
\multlinegap0pt
\usepackage{bm}
\usepackage{bbm}
\usepackage{dsfont}
\usepackage{xspace}
\usepackage{cancel}
\usepackage{mathtools}
\usepackage{amsthm}
\makeatletter
\@ifundefined{dfn}{\newtheorem{dfn}{Dfn}}{}
\@ifundefined{lem}{\newtheorem{lem}[dfn]{Lemma}}{}
\@ifundefined{prop}{\newtheorem{prop}[dfn]{Proposition}}{}
\makeatother
\title[Local energy decay of solutions to the linearized compressib]{Local energy decay of solutions to the linearized compressible viscoelastic system around motionless state in an exterior domain}
\author{Yusuke Ishigaki, Takayuki Kobayashi}
\date{Source: arXiv:2505.23169v1 (CC BY 4.0)}
\begin{document}
\maketitle

\noindent\emph{Source and licence.} Local energy decay of solutions to the linearized compressible viscoelastic system around motionless state in an exterior domain. Version arXiv:2505.23169v1 (CC BY 4.0), distributed under the Creative Commons Attribution 4.0 International licence, as recorded in the arXiv licence metadata for this exact version and shown on the abstract page https://arxiv.org/abs/2505.23169v1. Full rights notice: licenses/arxiv-2505.23169-CC-BY-4.0.txt.\par\smallskip

\noindent\textit{Editorial scope.} Counted unit: the Prop whose source label is \texttt{R-Ex-Prop2} in the article (source0.tex lines 890--917, source label \texttt{R-Ex-Prop2}), delivered together with the article's own complete original proof of it (source0.tex lines 1154--1209, one \texttt{proof} environment of 56 lines). No mathematics is written, repaired or re-derived: statement and proof are the article's own text line for line. The proof is detached from the statement in the source: the article states the result at lines 890--917 and prints its proof at lines 1154--1209, and the proof environment's own header names the statement, so the association is the article's own. Required context retained uncounted: R-Ex-1 (lines 181--185); R-R3-prop2 (lines 483--528); R-R3-prop3 (lines 731--747); R-Ex-Prop1 (lines 866--886); pr-R-Ex-Prop2-1-1 (lines 1001--1003); pr-R-Ex-Prop2-3 (lines 1059--1064); R-Ex-lem1 (lines 1071--1073). Every definition, display and label of those lines is retained unchanged. Omitted from this package and not redistributed: 1--180, 186--482, 529--730, 748--865, 887--889, 918--1000, 1004--1058, 1065--1070, 1074--1153, 1210--1408. Nothing inside a retained excerpt is omitted or altered: every retained body is delivered exactly as the article's lines are written, so no omission receipt of any kind accompanies this package. Citations: the retained excerpts carry no citation command at all, so no bibliography entry of the article is transcribed and no reference is redistributed. Reference adaptation: none was needed and none was made; every reference inside the delivered unit resolves inside it, so no unresolved reference or citation is produced. Printed slips kept literal: no printed word, symbol, hypothesis or number of the counted statement or of its proof was altered. Layout and class: the article's own class is \texttt{article} with packages including graphicx, amsmath, amssymb, amsthm, bm, latexsym, cases, tikz; the wrapper is the lane's pinned \texttt{amsart} editorial wrapper at 10pt on A4, which restates the theorem-like environments the retained text needs and adds the article's own macro definitions that the retained text needs (none), with extra wrapper packages (bm, bbm, dsfont, xspace, cancel, mathtools). The delivered numbering is the unit's own. Assets: no retained span contains a figure environment, a graphics-inclusion command, a PDF-page inclusion, a table or a TikZ picture; no image file is copied into this package, the package declares no asset dependency and no image file is redistributed with it. Licence: version arXiv:2505.23169v1 of this article is distributed under the Creative Commons Attribution 4.0 International licence, as recorded in the arXiv licence metadata for this exact version and shown on the abstract page https://arxiv.org/abs/2505.23169v1. A full rights notice is delivered at licenses/arxiv-2505.23169-CC-BY-4.0.txt. Characters: every retained excerpt is pure ASCII, so the delivered engine, XeLaTeX with its default Latin Modern text font, reports no missing character.
\begin{equation}
\label{R-Ex-1}
\lambda u+Au=\mathbf{f}~\text{in}~\Omega, ~
w|_{\partial \Omega}=0.
\end{equation}

\begin{prop}\label{R-R3-prop2}
Let  $1<p<\infty$ and $R, R^\prime>0$, and let $\lambda_1=\lambda_1(\nu,\tilde{\nu},\beta,\gamma)>0$ be a positive constant such that $\max\{|\lambda^2/\kappa_1(\lambda)|, |\lambda^2/\kappa_2(\lambda)|\}\le 1/2$ holds for any $\lambda\in U_{\lambda_1}$. Then $\mathcal{R}_{\mathbb{R}^3}^0(\lambda)$, $\Psi_{\mathbb{R}^3}^0(\lambda)$, $\mathcal{R}_{\mathbb{R}^3}^\infty(\lambda)$ and $\Psi_{\mathbb{R}^3}^\infty(\lambda)$ satisfy the following assertions.


\vskip2mm

\noindent
{\rm (i) (low frequency part)} Let $\alpha\in\mathbb{N}_0^3$, $\lambda\in \Dot{U}_{\lambda_1}$ and $\mathbf{f}\in L_R^p(\mathbb{R}^3)$. Then the two functions $\partial_x^\alpha \mathcal{R}_{\mathbb{R}^3}^0(\lambda)\mathbf{f}$ and $\partial_x^\alpha \Psi_{\mathbb{R}^3}^0(\lambda)\mathbf{f}$ are written as the sum of the singular and  holomorpic parts:
\begin{align*}
&\partial_x^\alpha \mathcal{R}_{\mathbb{R}^3}^0(\lambda)\mathbf{f}=\widetilde{\mathcal{R}}_{\mathbb{R}^3}^{0,\alpha}(\lambda)\mathbf{f}+\mathcal{R}_{\mathbb{R}^3}^{0,c,\alpha}(\lambda)\mathbf{f}, \quad
\partial_x^\alpha \Psi_{\mathbb{R}^3}^0(\lambda)\mathbf{f}=\widetilde{\Psi}_{\mathbb{R}^3}^{0,\alpha}(\lambda)\mathbf{f}+\Psi_{\mathbb{R}^3}^{0,c,\alpha}(\lambda)\mathbf{f}. 
\end{align*}
Here $\widetilde{\mathcal{R}}_{\mathbb{R}^3}^{0,\alpha}(\lambda)$ and $\widetilde{\Psi}_{\mathbb{R}^3}^{0,\alpha}(\lambda)$ are operator-valued functions satisfying 
\begin{align*}
&\widetilde{\mathcal{R}}_{\mathbb{R}^3}^{0,\alpha}(\lambda)\in \mathcal{A}(\dot{U}_{\lambda_1}\cap \Xi, \mathcal{L}(X_R^{p}(\mathbb{R}^3),L^{p}(B_{R^\prime})\times L^{p}(B_{R^\prime})^3 \times L^{p}(B_{R^\prime})^{3\times3} ) ),\\
&\widetilde{\Psi}_{\mathbb{R}^3}^{0,\alpha}(\lambda)\in \mathcal{A}(\dot{U}_{\lambda_1}\cap \Xi, \mathcal{L}(X_R^{p}(\mathbb{R}^3),L^{p}(B_{R^\prime})^3 ) )
\end{align*}
and the estimates
\begin{align*}
&\left\|\frac{d^m }{d\lambda^m}\widetilde{\mathcal{R}}_{\mathbb{R}^3}^{0,\alpha}(\lambda)\mathbf{f} \right\|_{L^p(B_{R^\prime})}\leq C\|\mathbf{f}\|_{X^p(\mathbb{R}^3)}
\begin{cases}
|\lambda|^{2-m} \left|\log |\lambda|\right| & (m\le 2), \\
|\lambda|^{2-m} &(m\ge3),
\end{cases}  
\\
& \left\| \frac{d^m}{d\lambda^m}\widetilde{\Psi}_{\mathbb{R}^3}^{0,\alpha} (\lambda)\mathbf{f}\right\|_{L^p( B_{R^\prime})} 
\le C\|\mathbf{f}\|_{X^p(\mathbb{R}^3)}
\begin{cases}
|\lambda|^{1-m}\left|\log |\lambda|\right| & (m\le 1), \\
|\lambda|^{1-m} &(m\ge2)
\end{cases}
\end{align*}
for $m\in\mathbb{N}_0, \lambda\in \Dot{U}_{\lambda_1}$ and $\mathbf{f}\in L_R^p(\mathbb{R}^3)$, where $C=C_{\alpha,m,p,R,R^\prime,\nu,\tilde{\nu},\beta,\gamma}$ is a positive constant; $\mathcal{R}_{\mathbb{R}^3}^{0,c,\alpha}(\lambda)$ and $\Psi_{\mathbb{R}^3}^{0,c,\alpha}(\lambda)$ are operator functions satisfying 
\begin{align*}
&\mathcal{R}_{\mathbb{R}^3}^{0,c,\alpha}(\lambda)\in \mathcal{A}(U_{\lambda_1} \cap \Xi, \mathcal{L}(X_R^{p}(\mathbb{R}^3),L^{p}(B_{R^\prime})\times L^{p}(B_{R^\prime})^3 \times L^{p}(B_{R^\prime})^{3\times3} )),\\
&\Psi_{\mathbb{R}^3}^{0,c,\alpha}(\lambda)\in \mathcal{A}(U_{\lambda_1} \cap \Xi, \mathcal{L}(X_R^{p}(\mathbb{R}^3),L^{p}(B_{R^\prime})^3)).
\end{align*}

\vskip2mm
\noindent
{\rm (ii) (high frequency part)} The operator-valued functions $\mathcal{R}_{\mathbb{R}^3}^\infty(\lambda)$ and $\Psi_{\mathbb{R}^3}^\infty(\lambda)$ satisfy 
\begin{align*}
&\mathcal{R}_{\mathbb{R}^3}^\infty(\lambda)\in \mathcal{A}(U_{\lambda_1}, \mathcal{L}(X^p(\mathbb{R}^3),Y^p(\mathbb{R}^3))),
~\Psi_{\mathbb{R}^3}^\infty(\lambda)\in \mathcal{A}(U_{\lambda_1},\mathcal{L}(X^p(\mathbb{R}^3),W^{2,p}(\mathbb{R}^3)^3) ).
\end{align*}
\end{prop}

\begin{prop}\label{R-R3-prop3}
Let $1<p<\infty$, $R,~R^\prime>0$ and $\mathbf{f}\in X_R^p(\mathbb{R}^3)$. Then $\mathcal{R}_{\mathbb{R}^3}(0)$ and $\Psi_{\mathbb{R}^3}(0)$ satisfies the following properties:
\begin{align}
&{\rm (i)}~\mathcal{R}_{\mathbb{R}^3}(0)\in \mathcal{L}(X_R^p(\mathbb{R}^3),Y^p(B_{R^\prime})),~\Psi_{\mathbb{R}^3}(0)\in \mathcal{L}(X_R^p(\mathbb{R}^3),W^{2,p}(B_{R^\prime})^3), \label{R-R3-prop3-1} \\
&{\rm (ii)}~ 
\left\|\nabla \Phi_{\mathbb{R}^3}(0)\mathbf{f}\right\|_{L^p(\mathbb{R}^3)}+\left\|\nabla \mathcal{W}_{\mathbb{R}^3}(0)\mathbf{f}\right\|_{W^{1,p}(\mathbb{R}^3)}+\left\|\nabla \mathcal{G}_{\mathbb{R}^3}(0)\mathbf{f}\right\|_{L^p(\mathbb{R}^3)}
\leq 
C\|\mathbf{f}\|_{X^p(\mathbb{R}^3)}, 
\label{R-R3-prop3-2}
\\
&{\rm (iii)}~
 |x|^2|(\mathcal{R}_{\mathbb{R}^3}(0)\mathbf{f})(x)|+|x||(\Psi_{\mathbb{R}^3}(0)\mathbf{f})(x)|\leq C \|\mathbf{f}\|_{X^p(\mathbb{R}^3)}~(|x|\ge 2R), 
\label{R-R3-prop3-3} \\
&{\rm (iv)}~
\lim_{|\lambda|\to 0,~|\arg \lambda|\le \frac{\pi}{2}}(\|(\mathcal{R}_{\mathbb{R}^3}(\lambda)-\mathcal{R}_{\mathbb{R}^3}(0))\mathbf{f}\|_{Y^p(B_{R^\prime})}+\|(\Psi_{\mathbb{R}^3}(\lambda)-\Psi_{\mathbb{R}^3}(0))\mathbf{f}\|_{W^{2,p}(B_{R^\prime})})=0. \label{R-R3-prop3-4}
\end{align}
\end{prop}

\begin{prop}\label{R-Ex-Prop1}
Let $1<p<\infty$, $0<\varepsilon_0<\pi/2$, $\lambda_0>0$ and $\mathbf{f}\in X^p(\Omega)$. 
Then, if $\lambda\in \Xi_{\varepsilon_0}$, then \eqref{R-Ex-1} has a unique solution $u=(\phi,w,G)=(\lambda+A)^{-1}\mathbf{f}$
satisfying the following properties.

\vskip3mm
\noindent
(i) The operator $(\lambda+A)^{-1}$ belongs to $\mathcal{A}(\Xi_{\varepsilon_0}; \mathcal{L}(X^p(\Omega); Y^p(\Omega) ) )$.

\vskip3mm
\noindent
(ii) The solution $u$ satisfies
\begin{equation}\label{R-Ex-prop1-1}
|\lambda | \|u\|_{X^p(\Omega)}+\sum_{j=1}^2|\lambda|^{\frac{1-j}{2}} \|\nabla^j w\|_{L^p(\Omega)}\leq C_{\varepsilon_0 ,\lambda_0}\|\mathbf{f}\|_{X^p(\Omega)}
\end{equation}
for $\lambda\in \Xi_{\varepsilon_0}$ with $|\lambda|\ge \lambda_0$. Additionally, $(\lambda+A)^{-1}$ also belongs to $\mathcal{A}(\Xi_{\varepsilon_0}; \mathcal{L}(Y^p(\Omega) ) )$ and $u$ is controlled as
\begin{equation}\label{R-Ex-prop1-2}
|\lambda | \|u\|_{Y^p(\Omega)}\leq C_{\varepsilon_0 ,\lambda_0}\|\mathbf{f}\|_{Y^p(\Omega)}
\end{equation}
for $\mathbf{f}\in Y^p(\Omega)$ and $\lambda\in \Xi_{\varepsilon_0}$ with $|\lambda|\ge \lambda_0$.
\end{prop}

\begin{align}
&\mathcal{S}^3(\lambda)\in \mathcal{A}(\dot{U}_{\lambda_1}, \mathcal{L}(X_b^p(\Omega)) ), \quad \|(\mathcal{S}^3(\lambda)-\mathcal{S}^3(0))\mathbf{f}\|_{X^p(\Omega_b)} \le C|\lambda|\|\mathbf{f}\|_{X^p(\Omega)}~(\lambda\in \dot{U}_{\lambda_1},~\mathbf{f}\in X_b^p(\Omega)). \label{pr-R-Ex-Prop2-1-1}
\end{align} 

\begin{gather}
(\lambda+A)\mathcal{R}^3(\lambda)\mathbf{f}=\mathbf{f}+\mathcal{S}^{3}(\lambda)\mathbf{f}~\textrm{in}~\Omega,~\mathcal{W}^3(\lambda)\mathbf{f}|_{\partial \Omega}=0, \label{pr-R-Ex-Prop2-3} \\
\int_{\Omega} \Phi^{3}(\lambda)\mathbf{f}dx=\int_{\Omega} \mathcal{S}_1^{3}(\lambda)\mathbf{f}dx=0, \quad \mathcal{G}^{3}(\lambda)\mathbf{f},~\mathcal{S}_1^{3}(\lambda)\mathbf{f} \in\nabla (W_0^{1,p}(\Omega))^3, \notag
\\
\nabla \Phi^3(\lambda)\mathbf{f}+\mathrm{div} {}^\top \mathcal{G}^3(\lambda)\mathbf{f}=\nabla \mathcal{S}_1^3(\lambda)\mathbf{f}+\mathrm{div} {}^\top \mathcal{S}_3^3(\lambda)\mathbf{f}=0~\mathrm{in}~\Omega,~\mathcal{S}^3(\lambda)\mathbf{f}=0~(|x|\ge b). \notag
\end{gather}

\begin{lem}\label{R-Ex-lem1}
For $1<p<\infty$, $I+\mathcal{S}^3(0)$ has its bounded inverse $(I+\mathcal{S}^3(0))^{-1}\in \mathcal{L}(X_b^p(\Omega))$.
\end{lem}

\begin{prop}\label{R-Ex-Prop2}
Let $1<p<\infty$, $0<\varepsilon_0<\pi/2$ and let $b, b_0>0$ be $\mathbb{R}^3 \setminus \Omega \subset B_{b_0}$ and $b>b_0$, and let $\mathbf{f}$ satisfy $\mathbf{f}\in X_b^p(\Omega)$. Let $D_{a_1,a_2}:=B_{a_2}\setminus B_{a_1}$ for $a_1,a_2\in \mathbb{R}$ with $a_2>a_1>0$.
Then, for any $b_1,b_2>0$ with $b_0<b_1<b_2<b$, there exist positive constant $\lambda_2=\lambda_2(\varepsilon_0,b_0,b_1,b_2,b,p)$ and operator $\widetilde{\mathcal{R}}_\Omega (\lambda)=\widetilde{\mathcal{R}}_{\varepsilon_0,b_0,b_1,b_2,b,p} (\lambda)$ such that the following assertions hold:
\begin{align}
{\rm (i)}~ &\widetilde{\mathcal{R}}_\Omega\in\mathcal{A}(\Dot{U}_{\lambda_2}; \mathcal{L}(X_b^p(\Omega),Y^p(\Omega_b))), \label{R-Ex-Prop2-1} \\
{\rm (ii)}~ &\widetilde{\mathcal{R}}_\Omega(\lambda)\mathbf{f}=(\lambda+A)^{-1}\mathbf{f}~(\lambda\in \Dot{U}_{\lambda_2}\cap \Xi_{\varepsilon_0},~\mathbf{f}\in X_b^p(\Omega) ), \label{R-Ex-Prop2-2} \\
{\rm (iii)}~
&\left\| \frac{d^m}{d\lambda^m}\widetilde{\mathcal{R}}_\Omega (\lambda)\mathbf{f}\right\|_{Y^p(\Omega_{b_1})} 
\le C\|\mathbf{f}\|_{X^p(\Omega)},
\label{R-Ex-Prop2-3} \\
&\left\| \frac{d^m}{d\lambda^m}\widetilde{\mathcal{R}}_\Omega (\lambda)\mathbf{f}\right\|_{Y^p( D_{b_1,b_2})} 
\le C\|\mathbf{f}\|_{X^p(\Omega)}
\begin{cases}
1 & (m=0), \\
\left|\log |\lambda|\right| & (m=1), \\
|\lambda|^{1-m} &(m\ge2),
\end{cases}  \label{R-Ex-Prop2-4} \\
&\left\| \frac{d^m}{d\lambda^m}\widetilde{\mathcal{R}}_\Omega (\lambda)\mathbf{f}\right\|_{Y^p(D_{b_2,b})} 
\le C\|\mathbf{f}\|_{X^p(\Omega)}
\begin{cases}
1 & (m\le1), \\
\left|\log |\lambda|\right| & (m=2), \\
|\lambda|^{2-m} &(m\ge3)
\end{cases}  \label{R-Ex-Prop2-5}
\end{align}
for $m\in\mathbb{N}_0$ and $\lambda\in \Dot{U}_{\lambda_2}$. 
Here $C=C_{\varepsilon_0,b_0,b_1,b_2,b,m,p}$ is a positive constant. 
\end{prop}

\begin{proof}[Proof of Proposition \ref{R-Ex-Prop2}]
We once assume the invertibility of $I+\mathcal{S}^3(\lambda)$ with small $\lambda$. Since $I+\mathcal{S}^3(\lambda)$ is rewritten as
$I+\mathcal{S}^3(\lambda)=(I+\mathcal{S}^3(0))\left[I+(I+\mathcal{S}^3(0))^{-1}(\mathcal{S}^3(\lambda)-\mathcal{S}^3(0))\right],$ we have the Neumann series:
\begin{equation}\label{pr-R-Ex-Prop2-4}
(I+\mathcal{S}^3(\lambda))^{-1}=\sum_{N=0}^\infty \left[(I+\mathcal{S}^3(0))^{-1} (\mathcal{S}^3(\lambda)-\mathcal{S}^3(0))\right]^N(I+\mathcal{S}^3(0))^{-1}
\end{equation}
in $\mathcal{L}(X_b^p(\Omega))$.
These facts are justified as the follows. The invertibility of $I+\mathcal{S}^3(0)$ is guaranteed by Lemma \ref{R-Ex-lem1}, and therefore we focus on the convergence of right-hand side of \eqref{pr-R-Ex-Prop2-4}. Noticing \eqref{pr-R-Ex-Prop2-1-1}, we can take $\lambda_2\in (0,\lambda_1]$ small such that
$
\left\| (I+\mathcal{S}^3(0))^{-1} (\mathcal{S}^3(\lambda)-\mathcal{S}^3(0))\right\|_{\mathcal{L}(X_b^p(\Omega))} \le 1/2$ for $\lambda\in \dot{U}_{\lambda_2}.
$
This implies that the right-hand side of \eqref{pr-R-Ex-Prop2-4} uniformly converges in $\mathcal{L}(X_b^p(\Omega))$. Consequently, $I+\mathcal{S}^3(\lambda)$ has its inverse for $\lambda\in \dot{U}_{\lambda_2}$ and $(I+\mathcal{S}^3(\lambda))^{-1}$ belongs to $ \mathcal{A}(\Dot{U}_{\lambda_2},\mathcal{L}(X_b^p(\Omega))).$
We then define $\widetilde{\mathcal{R}}_{\Omega}(\lambda)$ as
$
\widetilde{\mathcal{R}}_{\Omega}(\lambda):=\mathcal{R}^3(\lambda)(I+\mathcal{S}^3(\lambda))^{-1}.
$
Taking $(I+\mathcal{S}^3(\lambda))^{-1}$ in \eqref{pr-R-Ex-Prop2-3}, $\widetilde{\mathcal{R}}_{\Omega}(\lambda)\mathbf{f}$ solves \eqref{R-Ex-1} for $\lambda\in \dot{U}_{\lambda_2}$. Therefore, restricting $\lambda\in \dot{U}_{\lambda_2}\cap \Xi_{\varepsilon_0}$ and applying Proposition \ref{R-Ex-Prop1}, $\widetilde{\mathcal{R}}_{\Omega}(\lambda)\mathbf{f}=(\lambda+A)^{-1}\mathbf{f}$ is guranteed. This proves \eqref{R-Ex-Prop2-1} and \eqref{R-Ex-Prop2-2}. To show \eqref{R-Ex-Prop2-3} and \eqref{R-Ex-Prop2-4}, we first estimate the derivative of $\mathcal{R}^3(\lambda)\mathbf{f}$ with respect to $\lambda$. Let $m\ge 0$. Then the each components of $(d^m/d\lambda^m)\mathcal{R}^3(\lambda)\mathbf{f}$ are computed as
\begin{align*}
\frac{d^m}{d\lambda^m}\Phi^3(\lambda)\mathbf{f}&=(1-\varphi)\frac{d^m}{d\lambda^m}\Phi_{\mathbb{R}^3}(\lambda)\mathbf{f}_0+\frac{d^m}{d\lambda^m}(-\Delta_{\mathbb{R}^3})^{-1}\mathrm{div}\mathcal{G}_{\mathbb{R}^3}(\lambda)\mathbf{f}_0\cdot \nabla\varphi\\
&\quad +\delta_{m,0}\left\{\varphi\Phi_{\Omega_{b}}(\mathbf{f}|_{\Omega_{b}})+(-\Delta_{\Omega_b})^{-1}\mathrm{div}\mathcal{G}_{\Omega_b}(\mathbf{f}|_{\Omega_{b}})\cdot\nabla\varphi\right\}, \\
\frac{d^m}{d\lambda^m}\mathcal{W}^3(\lambda)\mathbf{f}&=(1-\varphi)\frac{d^m}{d\lambda^m} \mathcal{W}_{\mathbb{R}^3}(\lambda)\mathbf{f}_0+\delta_{m,0}\varphi\mathcal{W}_{\Omega_{b}}(\mathbf{f}|_{\Omega_{b}}), \\
\frac{d^m}{d\lambda^m}\mathcal{G}^3(\lambda)\mathbf{f}&=(1-\varphi)\frac{d^m}{d\lambda^m} \mathcal{G}_{\mathbb{R}^3}(\lambda)\mathbf{f}_0+ \frac{d^m}{d\lambda^m}(-\Delta_{\mathbb{R}^3})^{-1}\mathrm{div}\mathcal{G}_{\mathbb{R}^3}(\lambda)\mathbf{f}_0\otimes\nabla\varphi \\
&\quad + \delta_{m,0}\nabla\left\{\varphi(-\Delta_{\Omega_b})^{-1}\mathrm{div}\mathcal{G}_{\Omega_{b}}(\mathbf{f}|_{\Omega_{b}})\right\},
\end{align*}
where $\delta_{m,0}=1$ if $m=0$ and $\delta_{m,0}=0$ if $m\ge1$.
Using Proposition \ref{R-R3-prop2} and \ref{R-R3-prop3} to the resultant identities, we have
\begin{align*}
&\left\|\frac{d^m}{d\lambda^m}\mathcal{R}^3(\lambda)\mathbf{f}\right\|_{Y^p(\Omega_{b_1} )} 
\le C\|\mathbf{f}\|_{X^p(\Omega)}
\begin{cases}
1 & (m=0), \\
0 & (m\ge1),
\end{cases} \\
&\left\|\frac{d^m}{d\lambda^m}\mathcal{R}^3(\lambda)\mathbf{f}\right\|_{Y^p(D_{b_1,b_2})} 
\le C\|\mathbf{f}\|_{X^p(\Omega)}
\begin{cases}
1 & (m=0), \\
\left|\log |\lambda|\right| & (m=1), \\
|\lambda|^{1-m} &(m\ge2),
\end{cases} \\
&\left\|\frac{d^m}{d\lambda^m}\mathcal{R}^3(\lambda)\mathbf{f}\right\|_{Y^p( D_{b_2,b})} 
\le C\|\mathbf{f}\|_{X^p(\Omega)}
\begin{cases}
1 & (m\le 1), \\
\left|\log |\lambda|\right| & (m=2), \\
|\lambda|^{2-m} &(m\ge3)
\end{cases}
\end{align*}
for $\lambda\in \Dot{U}_{\lambda_2}$ and $\mathbf{f}\in X_b^p(\Omega)$.
Therefore, combining the above estimates, Leibniz' rule 
\[
\frac{d^m}{d\lambda^m}\widetilde{\mathcal{R}}(\lambda)\mathbf{f}=\sum_{l=0}^m \binom{m}{l} \frac{d^l}{d\lambda^l}\mathcal{R}^3(\lambda)\frac{d^{m-l}}{d\lambda^{m-l}}(I+\mathcal{S}^3(\lambda))^{-1}\mathbf{f}
\]
and $\|(d^m/d\lambda^m)(I+\mathcal{S}^3(\lambda))^{-1}\mathbf{f}\|_{X^p(\Omega)}\le C\|\mathbf{f}\|_{X^p(\Omega)}$ for $\lambda\in \Dot{U}_{\lambda_2}$ and $\mathbf{f}\in X_b^p(\Omega)$, we arrive at \eqref{R-Ex-Prop2-3}--\eqref{R-Ex-Prop2-5}.
This completes the proof.
\end{proof}

\end{document}
